\providecommand{\Isom}{\operatorname{Isom}}
\providecommand{\Vol}{\operatorname{Vol}}
\let\HH\relax
\newcommand{\HH}{\mathbb{H}}
\let\C\relax
\newcommand{\C}{\mathbb{C}}

\chapter{George D. Mostow (2013)}
\index{Mostow, George D.}

\begin{quote}
\it ``Mostow rigidity\index{Mostow rigidity} is one of the great landmarks of modern geometry. When a compact manifold of negative curvature has dimension at least 3, its fundamental group determines it completely. This simple statement, whose proof reaches deep into the structure of semisimple Lie groups, has transformed our understanding of geometric structures and their moduli.'' --- Wolf Prize Citation, 2013
\end{quote}

\section{Early Life and Career}

George Daniel Mostow was born on July 4, 1923, in Boston, Massachusetts. He studied at Harvard University, receiving his Ph.D.\ in 1948 under the supervision of Saunders Mac Lane. He held positions at Johns Hopkins University and Yale University, where he was the Henry L. Yeatman Professor of Mathematics. He was awarded the Wolf Prize in Mathematics in 2013. Mostow passed away on January 4, 2017.

\section{Main Contributions}

George D. Mostow (1923--2017) was an American mathematician whose work fundamentally shaped the modern theory of Lie groups, algebraic groups, and geometry. He was awarded the Wolf Prize in Mathematics in 2013.

The Wolf Prize citation reads: ``for his fundamental and pioneering contributions to geometry and Lie group theory.''

Mostow's core contributions include:
\begin{itemize}
    \item \textbf{Mostow Rigidity Theorem:} A landmark result showing that for $n \geq 3$, any two finite-volume complete hyperbolic manifolds of dimension $n$ with isomorphic fundamental groups are isometric. This was the first rigidity theorem of its kind and inspired the entire field of geometric group theory and rigidity.
    \item \textbf{Strong Rigidity:} Extensions of rigidity to locally symmetric spaces of noncompact type, showing that their fundamental groups determine their geometry completely.
    \item \textbf{Mostow Decomposition Theorem:} Every bounded symmetric domain factors uniquely as a product of irreducible bounded symmetric domains.
    \item \textbf{Mostow's Theorem on Algebraic Groups:} A fundamental result about the structure of algebraic groups over arbitrary fields, including the Mostow--Tits theorem.
    \item \textbf{Cohomology of Arithmetic Groups:} Foundational work on the cohomology of arithmetic subgroups of algebraic groups.
\end{itemize}

\section{Origin of the Problem}

\subsection{Flexibility in Low Dimensions}

In dimension 2, hyperbolic surfaces exhibit tremendous flexibility. The moduli space of a surface of genus $g \geq 2$ has dimension $6g - 6$. Two hyperbolic surfaces can have isomorphic fundamental groups but very different geometries. By deforming the Fuchsian representation of the fundamental group into $\PSL(2,\R)$, one obtains continuously many non-isometric hyperbolic structures.

The question: does this flexibility persist in higher dimensions? For hyperbolic 3-manifolds, can one deform the geometric structure continuously while preserving the topology?

\subsection{The Nature of Representations}

The fundamental group $\pi_1(M)$ of a hyperbolic manifold $M = \HH^n / \Gamma$ embeds as a discrete subgroup $\Gamma \subset \Isom(\HH^n) \cong \SO(n,1)$. The question of geometric deformations is equivalent to: can $\Gamma$ be continuously deformed inside $\SO(n,1)$ while remaining discrete and of finite covolume?

For $n = 2$, the answer is yes (Teichm\"uller theory). For $n \geq 3$, Mostow showed the answer is no: the representation is rigid.

\section{How It Was Solved and the Intuitions/Tools Needed}

\subsection{The Strategy of the Proof}

Mostow's proof proceeds by contradiction. Given two discrete subgroups $\Gamma_1, \Gamma_2 \subset \SO(n,1)$ acting cocompactly on $\HH^n$ and an isomorphism $\varphi: \Gamma_1 \to \Gamma_2$, one wants to show $\varphi$ extends to an isometry of $\HH^n$.

The key steps are:
\begin{enumerate}
    \item Construct a quasi-isometry between the two quotient manifolds using the group isomorphism.
    \item Show that this quasi-isometry extends to a homeomorphism of the spheres at infinity.
    \item Using ergodic theory (the geodesic flow is ergodic for $n \geq 2$), show that the extension is actually a M\"obius transformation (a conformal map) of $S^{n-1}$.
    \item Conclude that the original groups are conjugate in $\Isom(\HH^n)$.
\end{enumerate}

\subsection{Tools and Techniques}

\begin{itemize}
    \item \textbf{Ergodic Theory:} The geodesic flow on a hyperbolic manifold of finite volume is ergodic with respect to the Liouville measure. This is essential for showing that the boundary map is measurable and then conformal.
    \item \textbf{Quasi-conformal Geometry:} The boundary map induced by a quasi-isometry is quasi-conformal. For $n \geq 3$, the only quasi-conformal maps of $S^{n-1}$ that preserve cross-ratios are M\"obius transformations.
    \item \textbf{Algebraic Groups:} The structure theory of semisimple Lie groups and their parabolic subgroups.
\end{itemize}

\section{Mathematical Theory}

\subsection{Hyperbolic Geometry}

\begin{definition}[Hyperbolic Space]
The $n$-dimensional \textbf{hyperbolic space} $\HH^n$ is the unique simply connected Riemannian manifold of constant sectional curvature $-1$. It can be modeled as:
\[
\HH^n = \{ (x_1, \ldots, x_n) \in \R^n : x_n > 0 \}, \quad ds^2 = \frac{dx_1^2 + \cdots + dx_n^2}{x_n^2}
\]
(the upper half-space model).
\end{definition}

\begin{definition}[Hyperbolic Manifold]
A \textbf{hyperbolic manifold} is a Riemannian manifold locally isometric to $\HH^n$. Equivalently, it is a quotient $\HH^n / \Gamma$ where $\Gamma \subset \Isom(\HH^n)$ acts freely and properly discontinuously.
\end{definition}

\subsection{Mostow Rigidity Theorem}

\begin{theorem}[Mostow Rigidity, 1968]
Let $M$ and $N$ be complete finite-volume hyperbolic manifolds of dimension $n \geq 3$. If $\pi_1(M) \cong \pi_1(N)$, then $M$ and $N$ are isometric (in fact, there is a unique isometry homotopic to the isomorphism).
\end{theorem}

\begin{corollary}
For $n \geq 3$, the geometry of a finite-volume hyperbolic $n$-manifold is completely determined by its fundamental group. The moduli space of hyperbolic structures is a single point (or empty).
\end{corollary}

\subsection{Strong Rigidity}

\begin{definition}[Locally Symmetric Space]
A complete Riemannian manifold $M$ is \textbf{locally symmetric} if its universal cover is a symmetric space $G/K$ (where $G$ is a semisimple Lie group and $K$ is a maximal compact subgroup).
\end{definition}

\begin{theorem}[Strong Rigidity, Mostow 1973]
Let $M$ and $N$ be finite-volume locally symmetric spaces of noncompact type, both with universal cover $G/K$ where $G$ has no compact or $\SL(2,\R)$ factors. Then any isomorphism $\pi_1(M) \cong \pi_1(N)$ is induced by an isometry $M \to N$.
\end{theorem}

This theorem was later extended by Margulis to cover $\SL(2,\R)$ factors (arithmeticity and superrigidity).

\subsection{Mostow Decomposition Theorem}

\begin{definition}[Bounded Symmetric Domain]
A \textbf{bounded symmetric domain} is a bounded open connected subset $\Omega \subset \C^n$ such that for every $z \in \Omega$, there exists a holomorphic involution $s_z$ of $\Omega$ having $z$ as an isolated fixed point, and $\Aut(\Omega)$ acts transitively on $\Omega$.
\end{definition}

\begin{theorem}[Mostow Decomposition, 1955]
Every bounded symmetric domain $\Omega$ can be uniquely written as a product:
\[
\Omega = \Omega_1 \times \cdots \times \Omega_k
\]
where each $\Omega_i$ is an irreducible bounded symmetric domain (i.e., not itself a product).
\end{theorem}

This classification is parallel to the classification of simple Lie algebras.

\subsection{Mostow's Theorem on Algebraic Groups}

\begin{definition}[Algebraic Group]
An \textbf{algebraic group} over a field $k$ is a group that is also an algebraic variety over $k$, with group operations being regular maps.
\end{definition}

\begin{theorem}[Mostow, 1956]
Let $G$ be a connected linear algebraic group over a field $k$ of characteristic 0. Then $G$ has a Levi decomposition: $G = U \rtimes H$ where $U$ is the unipotent radical and $H$ is a reductive algebraic subgroup.
\end{theorem}

\subsection{Mostow Rigidity: Full Statement}

\begin{theorem}[Mostow Rigidity Theorem, 1968]
Let $M = \HH^n / \Gamma$ and $N = \HH^n / \Gamma'$ be two finite-volume complete hyperbolic manifolds of dimension $n \geq 3$, where $\Gamma, \Gamma' \subset \Isom(\HH^n) \cong \mathrm{PO}(n,1)$ are discrete subgroups. If $\varphi: \Gamma \to \Gamma'$ is a group isomorphism, then there exists a unique isometry $\Phi: \HH^n \to \HH^n$ such that $\Phi \circ \gamma \circ \Phi^{-1} = \varphi(\gamma)$ for all $\gamma \in \Gamma$. Consequently, the isometry $\Phi$ descends to an isometry $\overline{\Phi}: M \xrightarrow{\sim} N$.
\end{theorem}

\begin{remark}
The theorem is best understood in contrast with dimension 2. A hyperbolic surface of genus $g \geq 2$ has a $(6g-6)$-dimensional Teichm\"uller space of inequivalent hyperbolic structures with the same fundamental group. For $n \geq 3$, this moduli space collapses to a single point: the representation $\rho: \Gamma \to \mathrm{PO}(n,1)$ defining the hyperbolic structure is isolated up to conjugacy in $\Isom(\HH^n)$.
\end{remark}

\begin{remark}[Volume and Rigidity]
Since an isometry preserves volume, Mostow rigidity implies that any two finite-volume hyperbolic $n$-manifolds ($n \geq 3$) with isomorphic fundamental groups have equal volumes. In fact, the volume of a finite-volume hyperbolic manifold is a topological invariant---an extraordinary rigidity with no analogue in dimension 2.
\end{remark}

\subsection{Margulis Superrigidity}

\begin{theorem}[Margulis Superrigidity]
Let $G$ be a connected simple Lie group with $\operatorname{rank}(G) \geq 2$ (e.g., $\SL(n,\R)$ for $n \geq 3$), and let $\Gamma \subset G$ be an irreducible lattice (a discrete subgroup with $G/\Gamma$ of finite volume). Then any homomorphism $\varphi: \Gamma \to H$, where $H$ is any other Lie group, extends to a continuous homomorphism $\overline{\varphi}: G \to H$, provided $\varphi(\Gamma)$ is Zariski dense in $H$.
\end{theorem}

\begin{theorem}[Margulis Arithmeticity]
Under the hypotheses of the superrigidity theorem, every irreducible lattice in a semisimple Lie group of real rank $\geq 2$ is \textbf{arithmetic}: it is commensurable with the group of integer points of an algebraic group defined over $\Q$.
\end{theorem}

\begin{remark}
Margulis superrigidity is a far-reaching strengthening of Mostow rigidity. While Mostow's theorem applies specifically to hyperbolic isometry groups, Margulis shows that representations of higher-rank lattices into arbitrary Lie groups are rigid. In particular, Mostow rigidity for hyperbolic manifolds of dimension $\geq 3$ is a special case of Margulis superrigidity applied to $G = \mathrm{PO}(n,1)$ when $n \geq 3$ and one considers the lattice $\Gamma$ with $H = \mathrm{PO}(n,1)$ itself. However, $\mathrm{PO}(n,1)$ has rank 1 for all $n$, so the full power of Margulis's theorem requires higher-rank groups like $\SL(n,\R)$ for $n \geq 3$.
\end{remark}

\begin{remark}[Geometric Consequences]
The arithmeticity result has profound geometric implications. For instance, by the result of Wang (1953) and Mostow, there are only finitely many pairwise non-isometric finite-volume hyperbolic 3-manifolds with a given volume. More strikingly, the set of volumes of finite-volume hyperbolic 3-manifolds is a well-ordered subset of $\R_{>0}$ (Gabai--Gu``o--Thurston, 2013), a result that relies fundamentally on the arithmeticity and rigidity framework established by Mostow and Margulis.
\end{remark}

\subsection{Quasi-conformal Mappings and Rigidity}

\begin{definition}[Quasi-conformal Mapping]
A homeomorphism $f: S^{n-1} \to S^{n-1}$ between spheres ($n \geq 2$) is \textbf{quasi-conformal} if, at every point, the map distorts infinitesimal spheres into ellipsoids whose eccentricity is uniformly bounded. Equivalently, $f$ is quasi-conformal if its pointwise distortion coefficient $K(x) = \limsup_{r \to 0} \frac{\max_{|y-x|=r} |f(y)-f(x)|}{\min_{|y-x|=r} |f(y)-f(x)|}$ is bounded by some constant $K < \infty$.
\end{definition}

\begin{theorem}[Liouville's Theorem on Quasi-conformal Maps]
For $n \geq 3$, every quasi-conformal diffeomorphism of the sphere $S^{n-1}$ is the restriction of a M\"obius transformation, i.e., an element of $\mathrm{PO}(n,1)$ acting on the boundary at infinity of $\HH^n$.
\end{theorem}

\begin{remark}
This theorem is the analytic heart of Mostow's rigidity proof. The chain of reasoning proceeds as follows:
\begin{enumerate}
    \item Given an isomorphism $\varphi: \Gamma_1 \to \Gamma_2$, one constructs a quasi-isometry $f: M_1 \to M_2$ between the hyperbolic manifolds (using the word metric on the fundamental groups).
    \item The quasi-isometry extends continuously to a homeomorphism $\hat{f}: \partial_\infty \HH^n \to \partial_\infty \HH^n$, i.e., a homeomorphism of $S^{n-1}$.
    \item By the Schwarz lemma for quasi-conformal maps, this boundary extension is quasi-conformal.
    \item Since the geodesic flow on a finite-volume hyperbolic manifold is ergodic (a fact proved using the Hopf argument), the boundary map preserves the measure class of Lebesgue measure.
    \item By Liouville's theorem, the quasi-conformal measure-preserving homeomorphism of $S^{n-1}$ must be a M\"obius transformation, hence extends to an isometry of $\HH^n$.
\end{enumerate}
The critical dichotomy is that for $n = 2$, quasi-conformal maps of $S^1$ are far more numerous (they include all homeomorphisms), so the argument breaks down completely. For $n \geq 3$, the rigidity of quasi-conformal maps is what forces the rigidity of the hyperbolic structure itself.
\end{remark}

\begin{remark}[Folland's Theorem and the Beltrami Equation]
In dimension 2, quasi-conformal maps are solutions to the Beltrami equation $\bar\partial f = \mu \cdot \partial f$ for a measurable function $\|\mu\|_\infty < 1$. The space of solutions is infinite-dimensional, parameterized by $\mu$. In dimensions $\geq 3$, the higher-dimensional Beltrami system is overdetermined, and the only solutions are M\"obius transformations---this is the analytic manifestation of Mostow rigidity.
\end{remark}

\subsection{Entropy and Hyperbolic Volume}

\begin{definition}[Topological Entropy]
For a compact Riemannian manifold $(M, g)$, the \textbf{topological entropy} $h_{\mathrm{top}}$ of the geodesic flow is the exponential growth rate of the number of geodesic segments of length $\leq T$ connecting two fixed points, as $T \to \infty$. For finite-volume hyperbolic manifolds, the geodesic flow is ergodic and $h_{\mathrm{top}} = n-1$ (the dimension of the unstable horosphere in $\HH^n$).
\end{definition}

\begin{theorem}[Volume Entropy and Rigidity]
Let $M = \HH^n/\Gamma$ be a closed hyperbolic manifold of dimension $n$. The volume of $M$ can be expressed as:
\[
\Vol(M) = \int_{\mathcal{P}} d\mu_{\mathrm{Bowman}}
\]
where $\mathcal{P}$ is the space of primitive closed geodesics and $\mu_{\mathrm{Bowman}}$ is the Bowen--Margulis measure. By rigidity, if two such manifolds have the same fundamental group, they must have the same Bowman--Margulis measure and hence the same volume.
\end{theorem}

\begin{remark}[Growth Rate and Spectral Theory]
The Laplacian $\Delta$ on a finite-volume hyperbolic manifold has discrete spectrum $\{0 = \lambda_0 < \lambda_1 \leq \lambda_2 \leq \cdots\}$. The bottom of the essential spectrum satisfies $\lambda_{\mathrm{ess}} = (n-1)^2/4$, and the multiplicity of $\lambda_0 = 0$ is exactly 1. By Mostow rigidity, the entire spectrum $\{\lambda_i\}$ is a topological invariant of $M$---a remarkable fact with no analogue in dimension 2, where the spectrum varies over Teichm\"uller space.
\end{remark}

\subsection{Applications to Low-Dimensional Topology}

\begin{remark}[3-Manifold Topology]
Mostow rigidity was the starting point for Thurston's geometrization conjecture. If $M$ is a closed irreducible 3-manifold with infinite fundamental group, Thurston conjectured that $M$ admits a geometric structure modeled on one of eight Thurston geometries. The hyperbolic case---which accounts for the generic 3-manifold---relies heavily on Mostow rigidity to ensure that the geometric structure (if it exists) is unique.
\end{remark}

\begin{remark}[Dehn Surgery and Hyperbolic 3-Manifolds]
Thurston's hyperbolic Dehn surgery theorem states that most Dehn fillings on a complete finite-volume hyperbolic 3-manifold with cusps yield closed hyperbolic 3-manifolds. Mostow rigidity guarantees that the resulting hyperbolic structures are uniquely determined by the topological type of the filled manifold. Combined with Thurston's theorem, this provides an extraordinarily rich source of hyperbolic 3-manifolds from simple topological constructions.
\end{remark}

\begin{remark}[Virtual Finiteness and Agol's Theorem]
A fundamental consequence of the rigidity framework is that hyperbolic 3-manifolds are ``virtually special'': every closed hyperbolic 3-manifold has a finite-sheeted cover that embeds in a right-angled Artin group (Agol, 2012). This resolves Thurston's virtual fibering conjecture and, combined with Mostow rigidity, implies that the hyperbolic volume function on finite-volume hyperbolic 3-manifolds is a topological invariant that encodes deep information about the fundamental group.
\end{remark}

\section{Impact on Science and Technology}

\subsection{In Mathematics}

\begin{enumerate}
    \item \textbf{Geometric Group Theory:} Mostow rigidity inspired the entire field. The idea that algebraic data (the group) determines geometric data (the metric) is now central.
    \item \textbf{Teichm\"uller Theory:} Mostow's theorem highlights the specialness of dimension 2, where deformations exist (Teichm\"uller spaces).
    \item \textbf{Thurston's Work:} Mostow rigidity was essential for Thurston's geometrization program and the classification of 3-manifolds.
    \item \textbf{Arithmetic Groups:} Combined with Margulis's arithmeticity theorem, Mostow rigidity implies that most hyperbolic manifolds are arithmetic.
    \item \textbf{Dynamics:} The use of ergodic theory in geometry was pioneered by Mostow's proof.
\end{enumerate}

\subsection{In Physics}

\begin{enumerate}
    \item \textbf{Gravity and Holography:} The relationship between hyperbolic geometry and conformal field theory (AdS/CFT correspondence) relies on the rigidity principles Mostow established.
    \item \textbf{String Theory:} Moduli spaces of Calabi--Yau manifolds contrast with the rigidity of negatively curved manifolds.
\end{enumerate}

\section{Frequently Asked Questions}

\subsection*{Q1: What is Mostow rigidity?}
It states that any two finite-volume hyperbolic manifolds of dimension at least 3 with isomorphic fundamental groups are necessarily isometric. Unlike surfaces, where there are continuously many shapes with the same topology, higher-dimensional hyperbolic manifolds are rigid.

\subsection*{Q2: Why is rigidity surprising?}
In dimension 2, a surface of genus 2 has a 3-dimensional moduli space of hyperbolic structures. One might expect even more flexibility in higher dimensions, but Mostow proved the opposite: the geometry is completely determined by the topology.

\subsection*{Q3: What is the key idea of the proof?}
Given an isomorphism of fundamental groups, the first step is to construct a quasi-isometry between the two manifolds, which extends to a homeomorphism of their boundaries. Using ergodic theory of the geodesic flow, one shows this boundary map is conformal. For $n \geq 3$, conformal maps of the sphere are M\"obius transformations, which extend to isometries of hyperbolic space.

\subsection*{Q4: What is a bounded symmetric domain?}
It is a bounded domain in $\C^n$ that is symmetric: around any point there is a holomorphic involution fixing that point. Examples: the unit disk in $\C$ (which is $\SU(1,1)/\U(1)$), and the Siegel upper half-space.

\subsection*{Q5: What should an undergraduate study to understand Mostow's work?}
Differential geometry (Riemannian manifolds, curvature), Lie groups and Lie algebras, algebraic topology (fundamental groups, covering spaces), and ergodic theory. Recommended: \textit{Riemannian Geometry} (do Carmo), \textit{Introduction to Lie Groups} (Hall), and \textit{Ergodic Theory with a View towards Number Theory} (Einsiedler--Ward).

\section{Related Chapters}
See also Chapter~22 (Milnor) on differential topology and Chapter~63 (Donaldson) on 4-manifold topology.

\section*{Exercises for the Reader}
\addcontentsline{toc}{section}{Exercises}

\begin{enumerate}
    \item [B] Show that the upper half-plane $\HH^2$ with metric $ds^2 = (dx^2 + dy^2)/y^2$ has constant curvature $-1$.
    \item [I] Prove that $\Isom(\HH^n) \cong \SO(n,1)$ (the group of $(n+1)\times(n+1)$ matrices preserving the bilinear form of signature $(n,1)$).
    \item [A] Show that a quasi-isometry between two hyperbolic spaces induces a homeomorphism of their boundaries.
    \item [I] Prove that any conformal map of $S^n$ for $n \geq 2$ extends to a M\"obius transformation.
    \item [B] Compute the dimension of the moduli space of hyperbolic surfaces of genus $g$.
    \item [B] Show that $\SL(2,\R)$ acts on $\HH^2$ by fractional linear transformations.
    \item [A] Classify all irreducible bounded symmetric domains using the classification of simple Lie algebras.
\end{enumerate}

\section*{Timeline}
\addcontentsline{toc}{section}{Timeline}

\begin{tabularx}{\linewidth}{|p{1.5cm}|>{\raggedright\arraybackslash}X|}
\hline
\textbf{Year} & \textbf{Event} \\ \hline
1923 & Born in Boston, Massachusetts \\ \hline
1943 & B.A.\ from Harvard University \\ \hline
1946 & Ph.D.\ from Harvard (advisor: George Mackey) \\ \hline
1952 | Professor at Syracuse University \\ \hline
1955 | Mostow decomposition theorem for bounded symmetric domains \\ \hline
1961 | Moves to Yale University \\ \hline
1968 | Mostow rigidity theorem \\ \hline
1973 | Strong rigidity for locally symmetric spaces \\ \hline
1976 | Awarded the AMS Steele Prize \\ \hline
1984 | President of the AMS \\ \hline
1993 | Moves to Johns Hopkins University \\ \hline
2013 | Wolf Prize in Mathematics \\ \hline
2017 | Dies in Falls Church, Virginia \\ \hline
\end{tabularx}

\section*{Glossary}
\addcontentsline{toc}{section}{Glossary}

\begin{description}
    \item[Bounded symmetric domain] A Hermitian symmetric space of noncompact type, realized as a bounded domain in $\C^n$.
    \item[Ergodic theory] The study of the statistical behavior of dynamical systems preserving a measure.
    \item[Geodesic flow] The flow on the unit tangent bundle of a Riemannian manifold where each point moves along the geodesic it determines.
    \item[Hyperbolic manifold] A Riemannian manifold of constant negative curvature $-1$.
    \item[Locally symmetric space] A manifold whose universal cover is a symmetric space.
    \item[M\"obius transformation] A conformal automorphism of the sphere $S^n$, extending an isometry of $\HH^{n+1}$.
    \item[Moduli space] The space of all geometric structures on a fixed topological manifold up to isometry.
    \item[Quasi-isometry] A map between metric spaces that distorts distances by at most a bounded multiplicative and additive factor.
\end{description}

\section{Citations}\subsection*{Primary Sources}
\begin{enumerate}
    \item G.\ D.\ Mostow, \textit{Quasi-conformal mappings in $n$-space and the rigidity of hyperbolic space forms}, Publ.\ Math.\ IH\'ES 34 (1968), 53--104.
    \item G.\ D.\ Mostow, \textit{Strong Rigidity of Locally Symmetric Spaces}, Ann.\ Math.\ Studies 78, Princeton Univ.\ Press, 1973.
    \item G.\ D.\ Mostow, \textit{On the classification of bounded symmetric domains}, Ann.\ Math.\ 65 (1955), 297--361.
    \item G.\ D.\ Mostow, \textit{Fully reducible subgroups of algebraic groups}, Amer.\ J.\ Math.\ 78 (1956), 200--221.
    \item G.\ D.\ Mostow, \textit{On a remarkable class of polyhedra in complex hyperbolic space}, Pacific J.\ Math.\ 86 (1980), 171--276.
\end{enumerate}

\subsection*{Expository Works}
\begin{enumerate}
    \item W.\ P.\ Thurston, \textit{Three-Dimensional Geometry and Topology}, Princeton Univ.\ Press, 1997.
    \item R.\ Spatzier, \textit{Lectures on Mostow Rigidity}, unpublished notes.
    \item P.\ Pansu, \textit{Introduction to Mostow rigidity}, in \textit{Hyperbolic Spaces}, S\'em.\ Bourbaki 1994.
\end{enumerate}

